% =====================================================================
% Joint-Action Navigation: fairness statement + the two tables that matter.
% Every value was read back from the logged run configurations.
% =====================================================================
\subsection{Joint-Action Navigation: Protocol}
\label{app:jointnav-protocol}

\paragraph{Fair comparison.}
Within each task the four methods share everything except the actor. We
verified directly from the logged run configurations, rather than from the
launch scripts, that the environment, the reward terms, the PPO
hyperparameters, the exploration schedule, the off-policy correction and the
optimistic-advantage setting are identical across all four arms of every task.
Each task uses three seeds, and all four methods share the same seed set, the
same step budget and the same number of rollout workers. Crucially, the three
learning switches are \emph{not} specific to our method: group-commit
exploration, the arm-consistent off-policy weight and optimistic advantages are
applied to every arm, including \textsc{InforMARL}, whose training loop we
extended so that one command line configures both code bases.

\paragraph{Per-task settings are rule-based, not tuned.}
Two settings vary across tasks, each determined by a stated property of the
task rather than by sweeping. The subset-commit probability is $q{=}0.5$ on the
exact-$k$ tasks and $q{=}1$ otherwise: with $q{=}1$ an entire communication
group fires at once, which under an exact-$k$ rule is always an overkill, so
the three-fire/three-hold split could never be explored. Optimistic advantages
are enabled exactly when the miss penalty is $-5$, where exploration otherwise
yields almost uniformly negative advantages. Both settings apply to all four
methods (Table~\ref{tab:jointnav-settings}).

\begin{table}[h]
\centering
\caption{Task settings. Every row applies identically to all four methods.
All other hyperparameters are shared across tasks: $6$ agents, $3$ targets,
$k{=}3$, capture radius $0.3$, sensing/communication radius $0.6$, horizon
$50$, $\gamma{=}0.99$, $\lambda{=}0.95$, lr $5\times10^{-4}$, $5$ PPO epochs,
$32$ rollout workers, $3$ seeds.}
\label{tab:jointnav-settings}
\small
\begin{tabular}{lccccc}
\toprule
Task & Capture rule & Miss penalty & Exploration $\varepsilon$ & $q$ & Optimism \\
\midrule
E2     & at least $k$ & $-1$ & $0.2 \rightarrow 0$ & $1.0$ & -- \\
E2-p   & at least $k$ & $-5$ & $0.2$ constant      & $1.0$ & \checkmark \\
E4     & exactly $k$  & $-5$ & $0.2$ constant      & $0.5$ & \checkmark \\
E4-m1  & exactly $k$  & $-1$ & $0.2$ constant      & $0.5$ & -- \\
\bottomrule
\end{tabular}
\end{table}

\paragraph{What the results show.}
Three of the four tasks are solved by several methods, so the comparison does
not turn on whether a task can be solved at all. On E2, E2-p and E4-m1 the
methods that succeed reach comparable success rates, and the consistent
difference between them is \emph{efficiency}: \textsc{SDec-MAT} attains the
same outcome with roughly two to three times fewer unsuccessful in-range
capture attempts. Only the exact-$k$ task at the $-5$ penalty (E4) separates
the methods outright, and there every baseline fails while \textsc{SDec-MAT}
clears the board (Table~\ref{tab:jointnav-results}).

\begin{table}[h]
\centering
\caption{Joint-Action Navigation results, mean $\pm$ one standard deviation
over three seeds. \emph{False captures} is the number of unsuccessful in-range
capture attempts per episode (undershoots plus overkills; on the at-least-$k$
tasks no overkill is possible, so it equals the undershoot count). Lower is
better for false captures.}
\label{tab:jointnav-results}
\small
\begin{tabular}{llccc}
\toprule
Task & Method & Success rate & Joint captures & False captures $\downarrow$ \\
\midrule
\multirow{4}{*}{E2}
 & \textbf{SDec-MAT} & $0.88 \pm 0.02$ & $2.86 \pm 0.02$ & $\mathbf{0.49 \pm 0.05}$ \\
 & MAT-Dec           & $0.93 \pm 0.00$ & $2.92 \pm 0.01$ & $0.67 \pm 0.41$ \\
 & MAPPO             & $0.87 \pm 0.01$ & $2.78 \pm 0.04$ & $4.38 \pm 0.20$ \\
 & InforMARL         & $0.93 \pm 0.01$ & $2.88 \pm 0.00$ & $2.74 \pm 0.48$ \\
\midrule
\multirow{4}{*}{E2-p}
 & \textbf{SDec-MAT} & $0.88 \pm 0.01$ & $2.87 \pm 0.02$ & $\mathbf{0.79 \pm 0.21}$ \\
 & MAT-Dec           & $0.88 \pm 0.02$ & $2.86 \pm 0.02$ & $1.14 \pm 0.21$ \\
 & MAPPO             & $0.00 \pm 0.00$ & $0.72 \pm 0.43$ & $3.55 \pm 1.62$ \\
 & InforMARL         & $0.88 \pm 0.03$ & $2.87 \pm 0.03$ & $1.79 \pm 0.14$ \\
\midrule
\multirow{4}{*}{E4}
 & \textbf{SDec-MAT} & $\mathbf{0.65 \pm 0.39}$ & $\mathbf{2.47 \pm 0.60}$ & $1.97 \pm 1.33$ \\
 & MAT-Dec           & $0.00 \pm 0.00$ & $0.09 \pm 0.01$ & $1.42 \pm 0.11$ \\
 & MAPPO             & $0.00 \pm 0.00$ & $0.13 \pm 0.04$ & $2.18 \pm 0.21$ \\
 & InforMARL         & $0.00 \pm 0.00$ & $0.11 \pm 0.01$ & $1.72 \pm 0.13$ \\
\midrule
\multirow{4}{*}{E4-m1}
 & \textbf{SDec-MAT} & $0.85 \pm 0.08$ & $2.81 \pm 0.10$ & $\mathbf{1.80 \pm 0.89}$ \\
 & MAT-Dec           & $0.00 \pm 0.00$ & $0.12 \pm 0.01$ & $2.61 \pm 0.18$ \\
 & MAPPO             & $0.27 \pm 0.39$ & $1.00 \pm 1.24$ & $4.96 \pm 2.20$ \\
 & InforMARL         & $0.88 \pm 0.01$ & $2.83 \pm 0.03$ & $5.86 \pm 0.33$ \\
\bottomrule
\end{tabular}
\end{table}

\paragraph{Note.}
On E4 the baselines make few false captures only because they rarely attempt a
capture at all; per attempt, \textsc{SDec-MAT} succeeds $56\%$ of the time
against roughly $6\%$ for every baseline. Reported values are training-rollout
statistics, collected with exploration active, and therefore understate all
methods: greedy evaluation of the final E4 checkpoints over held-out layouts
gives $3.00$ captures at a $1.00$ success rate for \textsc{SDec-MAT}.
E2, E2-p and E4 use seeds $\{1,2,3\}$; E4-m1 was relaunched after a scheduler
fault and uses $\{2,3,4\}$, shared by all four methods.
Here's where things stand. The navigation table (tab:jointnav-results, Table 3 in the PDF) and the paper give me some of what the caption needs:

Seeds: 3, stated in Appendix F.3 and tab:jointnav-settings.
Training-rollout statistics with exploration active: already said in Appendix F.3 but missing from the Table 3 caption and the main-text discussion.

Two things aren't in the paper, and the public repo linked in the reproducibility statement is empty, so I can't confirm them yet:

What "±" is: standard deviation across the 3 seeds, standard error, or min–max. The MAPPO entry 27.0±39.0 suggests std across seeds, but I shouldn't guess.
The averaging window: for example, the last N rollout episodes per seed, or the final M training iterations.